\apendice{Documentación técnica de programación}

\section{Introducción}

En este apéndice se recoge la documentación técnica de programación del proyecto. Su objetivo es facilitar la comprensión de la estructura interna de la aplicación, así como proporcionar la información necesaria para que otro desarrollador pueda instalar, ejecutar, mantener o ampliar el sistema.

La aplicación desarrollada está formada por un \textit{backend} implementado con Django, un \textit{frontend} desarrollado con React y Vite, y una base de datos MySQL. El \textit{backend} se encarga de la lógica de negocio, la autenticación, la gestión de reservas, la carga de horarios y el acceso a la base de datos. El \textit{frontend}, por su parte, proporciona la interfaz de usuario desde la que se accede a las distintas funcionalidades del sistema.

Esta documentación describe la estructura de directorios del proyecto, las principales partes del código, los pasos necesarios para instalar y ejecutar la aplicación, y las pruebas realizadas para comprobar el correcto funcionamiento del sistema.

El proyecto se organiza separando el código del \textit{backend} y del \textit{frontend}, lo que permite mantener diferenciadas la lógica de negocio y la interfaz de usuario. Esta separación facilita el mantenimiento del sistema y permite que ambas partes puedan evolucionar de forma independiente.

\section{Estructura de directorios}

La estructura del proyecto se organiza separando la parte correspondiente al \textit{backend} y la parte correspondiente al \textit{frontend}. Esta división permite diferenciar claramente la lógica de negocio y acceso a datos, desarrollada con Django, de la interfaz de usuario, desarrollada con React. Por otro lado, tenemos el \textit{script SQL} de la base de datos llamado "BaseDeDatosHorario.sql", que se encuentra a la misma altura que el directorio del \textit{frontend} y del \textit{backend}.

En los siguientes apartados se describe la organización de ambos bloques, indicando la finalidad de los principales directorios y ficheros que forman parte del proyecto. Esta explicación facilita la comprensión del código fuente y permite que otros desarrolladores puedan localizar de forma más sencilla los módulos principales de la aplicación.


\subsection{Backend}

El \textit{backend} de la aplicación se encuentra desarrollado con Django y contiene la lógica principal del sistema. En esta parte del proyecto se gestionan los modelos de datos, las rutas, las vistas, la autenticación, la comunicación con la base de datos y las operaciones relacionadas con horarios, reservas, calendario académico, aulas, docencia y usuarios.

En la Figura~\ref{fig:estructura-directorios-backend} se muestra la estructura principal de directorios correspondiente al \textit{backend} del proyecto.

\imagen{img/EstructuraDirectoriosBack.png}{Estructura de directorios del backend}{fig:estructura-directorios-backend}

Como se observa en la Figura~\ref{fig:estructura-directorios-backend}, el directorio del \textit{backend} está formado por diferentes aplicaciones de Django, cada una de ellas orientada a una parte funcional concreta del sistema. Esta organización modular permite separar responsabilidades y facilita el mantenimiento del código.

El directorio \textit{nucleo} contiene la configuración principal del proyecto Django. En él se encuentran los archivos generales de configuración, como la configuración del proyecto, las rutas principales y otros elementos necesarios para el arranque de la aplicación.

El módulo \textit{aulas} agrupa la lógica relacionada con la gestión de aulas y su ocupación. Este módulo permite trabajar con la información relativa a capacidad, ubicación y equipamiento disponible en cada aula, y hacer las llamadas para obtener las reservas asociadas a cada aula.

El módulo \textit{calendario} contiene las funcionalidades relacionadas con el calendario académico. En él se gestiona la información de días lectivos, festivos, exámenes, cambios de docencia y fechas asociadas a TFG/TFM. Se encarga de cargar el calendario académico obteniendo los tipos de día y de manejar el cambio de tipo de día, además de la generación de un nuevo calendario académico.

El módulo \textit{docencia} agrupa la información académica relacionada con grados, asignaturas, grupos y docentes. Esta parte resulta fundamental para la generación y consulta de horarios, ya que permite asociar las reservas periódicas con la estructura docente correspondiente. En este módulo, se obtienen los cursos con horario cargado, los grados con docencia y sus semestres, las reservas periódicas asociadas a todo lo anterior, en cuanto a mostrar el horario. En cuanto a las reservas periódicas, maneja la creación de una nueva reserva periódica, la edición, su eliminación, las restricciones, la búsqueda de aulas disponibles frente a otras reservas periódicas.

El módulo \textit{app\_horarios} contiene las funcionalidades de importación y exportación de la base de datos conjunta o por entidades.

El módulo \textit{reservas} contiene la lógica relacionada con la gestión de reservas puntuales. En él se implementan las operaciones de creación, modificación, consulta, eliminación, y la carga de un horario creando con el analizador léxico y creando todas las reservas periódicas asociadas a la hoja de cálculo cargada.

El módulo \textit{usuarios} se encarga de las funcionalidades relacionadas con la autenticación, los usuarios y los roles. Esta parte permite controlar el acceso a las distintas funcionalidades de la aplicación en función del perfil asociado a cada usuario. Por lo que en esta parte se encuentra la conexión con la API del moodle de la universidad.

Además de estos directorios, aparecen varios ficheros relevantes en la raíz del proyecto. El archivo \textit{manage.py} es el punto de entrada para ejecutar comandos propios de Django, como iniciar el servidor de desarrollo, aplicar migraciones o crear usuarios administradores. El fichero \textit{requirements.txt} contiene las dependencias de Python necesarias para ejecutar el \textit{backend}. Por otro lado, los archivos \textit{package.json} y \textit{package-lock.json} recogen las dependencias asociadas a Node.js, utilizadas por la parte de \textit{frontend}. Finalmente, el directorio \textit{node\_modules} se genera automáticamente al instalar dichas dependencias.

A continuación, se adjunta la Tabla~\ref{tab:estructura-backend} a modo resumen de las funcionalidades que abarca cada módulo.

\begin{table}[H]
	\centering
	\begin{tabular}{l p{9cm}}
		\toprule
		\textbf{Directorio/Fichero} & \textbf{Descripción} \\
		\midrule
		\textit{nucleo} & Configuración principal del proyecto Django. \\
		\textit{aulas} & Gestión de aulas, capacidad, ubicación y recursos disponibles. \\
		\textit{calendario} & Gestión del calendario académico, días lectivos, festivos, exámenes y cambios de docencia. \\
		\textit{docencia} & Gestión de grados, asignaturas, grupos y docentes y horarios. \\
		\textit{app\_horarios} & Importación y exportación base de datos. \\
		\textit{reservas} & Gestión de reservas puntuales y analizador léxico. \\
		\textit{usuarios} & Gestión de autenticación, usuarios, roles y permisos. \\
		\textit{manage.py} & Archivo principal para ejecutar comandos de administración de Django. \\
		\textit{requirements.txt} & Fichero con las dependencias Python necesarias para el backend. \\
		\textit{package.json} & Fichero de configuración de dependencias de Node.js. \\
		\textit{node\_modules} & Directorio generado automáticamente al instalar dependencias de Node.js. \\
		\bottomrule
	\end{tabular}
	\caption{Descripción de la estructura principal del backend}
	\label{tab:estructura-backend}
\end{table}

\subsection{Frontend}

El \textit{frontend} de la aplicación se encuentra desarrollado con React + Vite y contiene la parte visual del sistema. Esta parte del proyecto se encarga de mostrar las pantallas de la aplicación, gestionar la interacción con el usuario y comunicarse con el \textit{backend} mediante peticiones HTTP.

En la Figura~\ref{fig:estructura-directorios-frontend} se muestra la estructura principal de directorios correspondiente al \textit{frontend} del proyecto.

\imagen{img/EstructuraDirectoriosFront.png}{Estructura de directorios del frontend}{fig:estructura-directorios-frontend}

Como se observa en la Figura~\ref{fig:estructura-directorios-frontend}, el directorio del \textit{frontend} está formado por una carpeta principal \textit{src}, donde se encuentra el código fuente de la interfaz, y por varios ficheros de configuración necesarios para ejecutar y construir la aplicación.

El directorio \textit{public} contiene recursos estáticos públicos que pueden ser utilizados directamente por la aplicación. Estos archivos se sirven sin necesidad de ser procesados por el sistema de construcción.

Dentro de \textit{src} se organiza la mayor parte del código del \textit{frontend}.

El directorio \textit{administrador} agrupa las pantallas y componentes relacionados con el panel de administración. En este módulo se encuentran las funcionalidades destinadas a la gestión de entidades, usuarios, roles y el mantenimiento global del sistema.

El directorio \textit{api} contiene las funciones encargadas de centralizar la comunicación con el \textit{backend}. Esta organización permite separar las llamadas HTTP de los componentes visuales, facilitando el mantenimiento del código y evitando duplicidades.

El directorio \textit{assets} almacena recursos estáticos utilizados por la interfaz, como imágenes, logotipos u otros elementos gráficos.

El directorio \textit{aulas} agrupa los componentes y páginas relacionados con la consulta de ocupación de aulas.

El directorio \textit{auth} contiene la lógica para proteger rutas de acceso sin autenticación, y para validar permisos.

El directorio \textit{calendario} agrupa las pantallas y componentes relacionados con la consulta y modificación del calendario académico.

El directorio \textit{horarios} contiene las pantallas y componentes necesarios para cargar nuevos horarios, consultar horarios existentes.

El directorio \textit{login} contiene la pantalla de inicio de sesión y los elementos visuales relacionados con el acceso a la aplicación.

El directorio \textit{reservas} agrupa las pantallas y componentes relacionados con la gestión de reservas puntuales y periódicas, incluyendo listados, formularios de creación y edición, y acciones asociadas a las solicitudes.

Finalmente, el directorio \textit{shared} contiene elementos compartidos por diferentes módulos del \textit{frontend}, permitiendo reutilizar código común y mantener una estructura más organizada.

Además de estos directorios, aparecen varios ficheros relevantes en la raíz del proyecto. El archivo \textit{App.jsx} define la estructura principal de la aplicación React y la organización general de sus rutas o componentes principales. El archivo \textit{main.jsx} actúa como punto de entrada de la aplicación. Los ficheros \textit{App.css} e \textit{index.css} contienen estilos globales utilizados por la interfaz.

El fichero \textit{package.json} recoge la configuración del proyecto, los scripts de ejecución y las dependencias necesarias. El archivo \textit{package-lock.json} almacena la versión exacta de las dependencias instaladas. El archivo \textit{vite.config.js} contiene la configuración de Vite, herramienta utilizada para ejecutar y construir el \textit{frontend}. Por último, el directorio \textit{node\_modules} se genera automáticamente al instalar las dependencias de Node.js y no forma parte del código fuente desarrollado.

Se ha realizado una Tabla~\ref{tab:estructura-frontend} a modo resumen visual de lo que realiza cada directorio.

\begin{table}[H]
	\centering
	\begin{tabular}{l p{9cm}}
		\toprule
		\textbf{Directorio/Fichero} & \textbf{Descripción} \\
		\midrule
		\textit{public} & Recursos estáticos públicos de la aplicación. \\
		\textit{src} & Directorio principal del código fuente del frontend. \\
		\textit{administrador} & Pantallas y componentes relacionados con el panel de administración. \\
		\textit{api} & Funciones encargadas de realizar peticiones al backend. \\
		\textit{assets} & Imágenes, logotipos y otros recursos gráficos. \\
		\textit{aulas} & Componentes relacionados con aulas y ocupación de espacios. \\
		\textit{auth} & Lógica de autenticación y gestión de sesión. \\
		\textit{calendario} & Componentes relacionados con el calendario académico. \\
		\textit{general} & Componentes comunes de la interfaz. \\
		\textit{horarios} & Pantallas de carga, consulta y gestión de horarios. \\
		\textit{login} & Pantalla y componentes de inicio de sesión. \\
		\textit{reservas} & Gestión de reservas puntuales. \\
		\textit{shared} & Componentes o utilidades compartidas por varios módulos. \\
		\textit{App.jsx} & Componente principal de la aplicación React. \\
		\textit{main.jsx} & Punto de entrada del frontend. \\
		\textit{vite.config.js} & Configuración de Vite. \\
		\textit{package.json} & Dependencias y scripts del proyecto frontend. \\
		\textit{node\_modules} & Dependencias instaladas automáticamente mediante Node.js. \\
		\bottomrule
	\end{tabular}
	\caption{Descripción de la estructura principal del frontend}
	\label{tab:estructura-frontend}
\end{table}


\section{Manual del programador}

Esta sección está dirigida a futuros desarrolladores que necesiten mantener, ampliar o modificar la aplicación. En ella se describen los aspectos más relevantes del desarrollo, la organización interna del código y las partes principales que deben conocerse antes de realizar cambios en el sistema.

La aplicación está formada por dos bloques principales: el \textit{backend}, desarrollado con Django, y el \textit{frontend}, desarrollado con React. El \textit{backend} se encarga de la lógica de negocio, la autenticación, la comunicación con la base de datos y la validación de las operaciones. El \textit{frontend}, por su parte, proporciona la interfaz visual y se comunica con el servidor mediante peticiones HTTP.

\subsection{Organización interna del backend}

Como se ha indicado en la sección anterior, el \textit{backend} se encuentra dividido en varios módulos funcionales, como \textit{aulas}, \textit{calendario}, \textit{docencia}, \textit{reservas} y \textit{usuarios}. Cada uno de estos módulos agrupa la lógica correspondiente a una parte concreta del sistema, lo que permite separar responsabilidades y facilita el mantenimiento del código.

Dentro de estos módulos se ha seguido una organización basada en distintos tipos de ficheros, entre los que destacan los modelos, vistas, serializadores y servicios.

\subsubsection{Modelos}

Los modelos representan las entidades de la base de datos dentro de Django. Cada modelo define los campos de una tabla, sus tipos de datos, relaciones y restricciones principales. Gracias a estos modelos, el \textit{backend} puede trabajar con la información persistente de la aplicación sin necesidad de escribir directamente todas las consultas SQL.

En este proyecto, los modelos representan elementos como aulas, docentes, asignaturas, grupos, reservas, cursos, semestres y días del calendario académico. Cualquier modificación en la estructura de datos debe revisarse cuidadosamente, ya que puede afectar a las vistas, serializadores, servicios y pantallas del \textit{frontend} que dependan de dicha información.

\subsubsection{Vistas}

Las vistas son las encargadas de recibir las peticiones procedentes del \textit{frontend}, procesarlas y devolver una respuesta. En ellas se definen los puntos de entrada de la API y se controla qué operación debe ejecutarse en función de la petición recibida.

El \textit{frontend} envía una petición HTTP al \textit{backend}. La vista correspondiente recibe esa petición, comprueba los datos recibidos y coordina el resto de operaciones necesarias.

Aunque las vistas pueden contener validaciones sencillas, se ha procurado que la lógica más compleja se separe en servicios auxiliares. De esta forma, las vistas se mantienen más claras y se evita concentrar demasiada lógica en un único fichero.

\subsubsection{Serializadores}

Los serializadores se utilizan para transformar los datos entre el formato interno de Django y el formato intercambiado con el \textit{frontend}, normalmente JSON. Su función es especialmente importante en una aplicación con \textit{frontend} y \textit{backend} desacoplados, ya que permiten controlar qué información se envía al cliente y cómo se reciben los datos introducidos por el usuario.

Además de convertir objetos en respuestas JSON, los serializadores permiten realizar validaciones sobre los datos recibidos antes de crear o actualizar registros en la base de datos. Por ejemplo, pueden comprobar que los campos obligatorios estén presentes, que los tipos de datos sean correctos o que una estructura enviada desde el \textit{frontend} cumpla el formato esperado.

\subsubsection{Servicios}

Los servicios agrupan funciones auxiliares o procesos de negocio que no deben estar directamente dentro de las vistas. Esta separación permite reutilizar código y mantener una estructura más organizada.

En este proyecto, los servicios son especialmente importantes en procesos como la carga de horarios desde hojas de cálculo, la validación de disponibilidad de aulas, la comprobación de solapes entre reservas y el tratamiento del calendario académico.

Por ejemplo, en lugar de implementar toda la lógica de lectura de Excel dentro de una vista, esta se puede delegar en funciones o clases de servicio encargadas de procesar el fichero, extraer la información relevante y devolver los datos estructurados para que posteriormente sean validados y almacenados.

\subsubsection{Rutas}

Las rutas definen las direcciones de la API disponibles para el \textit{frontend}. Cada ruta se asocia con una vista o conjunto de vistas que implementan una funcionalidad concreta. Mantener las rutas organizadas por módulos facilita localizar los puntos de entrada de cada parte del sistema.

Si en el futuro se añade una nueva funcionalidad, será necesario definir su ruta correspondiente, implementar la vista asociada, crear o reutilizar los serializadores necesarios y, en caso de que exista lógica compleja, trasladarla a un servicio específico.

\subsection{Organización interna del frontend}

El \textit{frontend} está desarrollado con React y se organiza en módulos funcionales relacionados con las principales partes de la aplicación, como reservas, horarios, calendario, aulas, autenticación y administración.

Cada módulo contiene componentes y pantallas relacionados con una funcionalidad concreta. Esta organización permite que el código sea más mantenible y que los cambios realizados en una parte de la interfaz no afecten directamente al resto del sistema.

\subsubsection{Componentes}

Los componentes son la base de la interfaz en React. Cada componente representa una parte visual o funcional de la aplicación, como formularios, tablas, botones, modales, menús o pantallas completas.

En este proyecto se han utilizado componentes para construir pantallas como la gestión de reservas puntuales, la ocupación de aulas, el calendario académico, la carga de horarios o el panel de administración. La reutilización de componentes permite evitar duplicidades y mantener una apariencia coherente en toda la aplicación.

\subsubsection{Estado de los componentes}

Para gestionar la información interna de los componentes se ha utilizado el \textit{hook} \textit{useState}. Este mecanismo permite almacenar valores que pueden cambiar durante la ejecución de la aplicación, como filtros seleccionados, datos de formularios, mensajes de error, resultados de una consulta o el estado de apertura de un modal.

\subsubsection{Efectos y carga de datos}

Para ejecutar acciones asociadas al ciclo de vida de los componentes se ha utilizado el \textit{hook} \textit{useEffect}. Este mecanismo permite realizar operaciones cuando se carga una pantalla o cuando cambia algún dato del que depende el componente.

En este proyecto, \textit{useEffect} se utiliza principalmente para cargar información inicial desde el \textit{backend}. También puede emplearse para actualizar la información mostrada cuando el usuario modifica filtros o selecciona una nueva opción.

De esta forma, las pantallas pueden solicitar datos al \textit{backend} en el momento adecuado y actualizar la interfaz cuando la respuesta ha sido recibida.

\subsubsection{Comunicación con el backend}

La comunicación entre el \textit{frontend} y el \textit{backend} se realiza mediante peticiones HTTP. Para evitar duplicar código, las llamadas a la API se centralizan en ficheros específicos dentro del directorio \textit{api}.

Estos ficheros contienen funciones encargadas de realizar operaciones concretas, como consultar reservas, crear una reserva puntual, modificar una reserva periódica, cargar un horario o actualizar un día del calendario académico. Los componentes de React llaman a estas funciones cuando necesitan comunicarse con el servidor.

Esta organización permite separar la lógica visual de la lógica de comunicación. Si en el futuro cambia una ruta del \textit{backend}, solo será necesario actualizar el fichero correspondiente de la API, sin modificar todas las pantallas que utilizan esa operación.

\subsubsection{Gestión de formularios y eventos}

Las pantallas que permiten crear o editar información utilizan formularios controlados mediante estado. Cada campo del formulario se asocia a una variable de estado y se actualiza mediante eventos, como \textit{onChange} o \textit{onSubmit}.

Cuando el usuario confirma una operación, el componente recoge los valores almacenados en el estado, realiza las comprobaciones necesarias en la interfaz y envía los datos al \textit{backend}. Posteriormente, en función de la respuesta recibida, se muestra un mensaje de confirmación o de error.

\subsection{Control de autenticación y permisos}

El sistema diferencia las funcionalidades disponibles en función del rol asociado al usuario autenticado. Por este motivo, el programador debe tener en cuenta que algunas pantallas, botones o acciones solo deben mostrarse a determinados perfiles.

En el \textit{frontend}, esta restricción se utiliza para adaptar la interfaz al rol del usuario, ocultando opciones que no le corresponden. Sin embargo, estas comprobaciones no son suficientes por sí solas, ya que un usuario podría intentar acceder manualmente a una ruta o enviar una petición no permitida.

Por este motivo, las validaciones importantes deben realizarse también en el \textit{backend}, que es el encargado de garantizar que cada operación solo pueda ser ejecutada por usuarios autorizados.


\section{Compilación, instalación y ejecución del proyecto}

\section{Pruebas del sistema}
